% ============================================================================
\uxsection{Introducción}

Esta actividad da continuidad a la definición de producto y de usuarios que el equipo
presentó en la actividad anterior. Allí se caracterizó la audiencia de \textbf{Karisma Data},
el portal centralizado de datos financieros que propone el equipo, y se desarrollaron ocho
personas con sus respectivos mapas de empatía. Aquellos artefactos describen quién es cada
usuario, qué necesita y qué siente, pero lo hacen de forma estática. El objetivo de esta
actividad es distinto: definir la interacción del usuario con el producto \emph{a través del
tiempo}.

Para lograrlo se produjeron dos artefactos. El primero es un conjunto de \textbf{seis
escenarios}, dos por integrante del equipo, elaborados de forma individual a partir de las
personas y los mapas de empatía de la actividad anterior. Cada escenario sitúa a una persona
concreta en un contexto real de trabajo, describe la tarea que quiere completar, los pasos
que sigue, los puntos en los que entra en contacto con el producto y el resultado que espera
obtener, incluida la sensación que busca al conseguirlo. El segundo artefacto son
\textbf{cuatro journey maps}: tres elaborados de forma individual, uno por integrante sobre
uno de sus propios escenarios, y el \textbf{journey map de equipo}, construido en un taller
conjunto. Cada mapa toma el objetivo de un escenario, lo descompone en un listado de
actividades y lo despliega en una línea de tiempo con las etapas, las acciones, los puntos de
contacto, los pensamientos, las emociones, los puntos de fricción y las oportunidades de
diseño que se derivan de ellos. Los cuatro comparten la misma estructura de siete carriles y
la misma escala emocional, de modo que puedan compararse entre sí.

El método de trabajo fue el siguiente. Cada integrante escribió los escenarios de las personas
que había desarrollado en la actividad anterior, siguiendo una plantilla común que se describe
en la sección~\ref{sec:metodo} y que proviene del material de apoyo del curso, y construyó
después el mapa de uno de ellos. El equipo discutió los seis escenarios en conjunto,
seleccionó el objetivo de la persona primaria del producto como base del mapa de equipo y lo
elaboró en un taller de dos horas cuyo diseño se documenta en el apartado~\ref{sec:taller}.
Los escenarios y los mapas constituyen, en esta etapa, hipótesis de trabajo fundamentadas en
la caracterización de la audiencia, en la experiencia profesional del equipo en áreas que
operan con información financiera institucional y en la literatura de diseño centrado en el
usuario; el apartado~\ref{sec:alcance} declara con precisión ese alcance y el plan previsto
para validarlo con evidencia de campo.

El documento se organiza en tres bloques. El primero explica el método, la plantilla de
escenario empleada y el alcance declarado. El segundo reúne el trabajo individual: los seis
escenarios y los tres journey maps elaborados por cada integrante, agrupados bajo su nombre.
El tercero desarrolla el journey map de equipo y conecta sus hallazgos con el producto,
mediante una matriz que traza cada punto de fricción hasta la característica que lo resuelve y
la métrica que lo verificará, y un conjunto de principios rectores y antiprincipios derivados
del ejercicio. El Anexo~A reproduce el journey map de equipo en formato de tabla.
